\section{Task Topologies and Representation Requirements}
\label{sec:topologies}

\subsection{Why topology is a more useful organizing variable than task name}

Task labels conceal important physical differences. ``Handover'' can mean a
quasi-static transfer with simultaneous support, a dynamic throw--catch event,
or a robot--human release triggered by pull and slip. ``Tool use'' can mean one
hand moving a rigid tool against a fixed environment, one hand stabilizing a
workpiece while the other acts, or both hands jointly holding the same tool.
These variants demand different information even when their semantic label is
the same. A representation-oriented review should therefore begin with the
active interfaces and their changes.

Let the task at time $t$ contain entities
$\mathcal{V}_t=\{H_L,H_R,O_1,\ldots,T_1,\ldots,W,E\}$ for left and right hands,
objects, tools, a workpiece, and environment. The physical contact topology is
$\mathcal{G}^{c}_t=(\mathcal{V}_t,\mathcal{E}^{c}_t)$, where an edge represents
an active load-bearing or constraint-producing interface. This notation does
not require the model itself to be a graph neural network. It provides a neutral
way to ask whether the observation history $o_{1:t}$ contains enough information
to infer the relevant interfaces and their direction-dependent effects.

\begin{figure}[t]
\centering
\begin{tikzpicture}[
  entity/.style={circle, draw=reviewblue, fill=reviewlight, minimum size=9mm,
                 font=\small\bfseries},
  tool/.style={rectangle, rounded corners, draw=revieworange, fill=reviewwarm,
               minimum width=12mm, minimum height=8mm, font=\small\bfseries},
  contact/.style={line width=1.2pt, draw=reviewteal},
  inactive/.style={line width=0.8pt, draw=reviewgray, dashed},
  phase/.style={font=\small\bfseries, text=reviewgray}
]
% handover phases
\node[phase] at (1.8,2.65) {Quasi-static handover};
\node[entity] (gl) at (0,1.8) {$H_L$};
\node[tool] (o1) at (1.6,1.8) {$O$};
\node[entity] (gr) at (3.2,1.8) {$H_R$};
\draw[contact] (gl)--(o1);
\draw[inactive] (o1)--(gr);
\node[font=\scriptsize] at (1.6,1.25) {giver-only};

\node[entity] (dl) at (4.3,1.8) {$H_L$};
\node[tool] (o2) at (5.9,1.8) {$O$};
\node[entity] (dr) at (7.5,1.8) {$H_R$};
\draw[contact] (dl)--(o2);
\draw[contact] (o2)--(dr);
\node[font=\scriptsize] at (5.9,1.25) {dual contact};

\node[entity] (rl) at (8.6,1.8) {$H_L$};
\node[tool] (o3) at (10.2,1.8) {$O$};
\node[entity] (rr) at (11.8,1.8) {$H_R$};
\draw[inactive] (rl)--(o3);
\draw[contact] (o3)--(rr);
\node[font=\scriptsize] at (10.2,1.25) {receiver-only};

% tool workpiece
\node[phase] at (6.0,0.35) {Cooperative tool--workpiece manipulation};
\node[entity] (hl2) at (2.0,-0.65) {$H_L$};
\node[tool] (tool2) at (4.2,-0.65) {$T$};
\node[tool] (work) at (7.3,-0.65) {$W$};
\node[entity] (hr2) at (9.5,-0.65) {$H_R$};
\draw[contact] (hl2)-- node[above,font=\scriptsize] {grasp} (tool2);
\draw[contact] (tool2)-- node[above,font=\scriptsize] {working contact} (work);
\draw[contact] (work)-- node[above,font=\scriptsize] {support} (hr2);
\node[tool] (env) at (7.3,-1.75) {$E$};
\draw[inactive] (work)-- node[right,font=\scriptsize] {optional} (env);
\end{tikzpicture}
\caption{Two complementary topology families. Handover changes which hand
supports a shared object; tool--workpiece manipulation introduces a longer,
role-asymmetric load path. Dashed edges are inactive or optional.}
\label{fig:topologies}
\end{figure}

\subsection{Six topology families}

\paragraph{Local single-interface contact.}
The simplest case associates a tactile patch with local geometry, force, slip,
or material. This is the dominant pretraining unit for image-based tactile
encoders. It supports transferable perception but leaves global interface
identity and cross-body coupling to downstream models.

\paragraph{Multi-finger contact within one hand.}
In-hand manipulation requires several contacts to be fused across finger
kinematics and time. Distributed-taxel graphs and hierarchical tactile graphs
explicitly exploit this structure~\cite{funabashi2022distributed,yang2023tacgnn}.
The object couples the fingers, making this topology a useful precursor to the
two-hand problem, but it does not include inter-hand role changes or object
transfer.

\paragraph{Two hands sharing one object.}
Both hands may grasp, reorient, stretch, or stabilize the same object. Bimanual
simulation suites and real policy-learning systems cover many such
tasks~\cite{chen2022bidexhands,shaw2024bidex,jiang2024dexmimicgen}.
Success can arise from kinematic coordination and demonstration priors even when
load sharing is not explicitly represented.

\paragraph{Handover.}
Handover is a topology transition rather than merely a motion trajectory. Static
and quasi-static transfer offers giver-only, dual-contact, and receiver-only
segments; dynamic throw--catch adds a flight phase and impact. The latter is an
impressive control test~\cite{huang2023dynamichandover} but a weaker primary
benchmark for continuous cross-hand load propagation because the hands are not
simultaneously connected during flight. Controlled overlapping handover is
better suited to studying support redistribution, release readiness, and
transition uncertainty~\cite{costanzo2021handover,wang2024contacthandover}.

\paragraph{One hand, tool, and environment.}
Tool contact creates extrinsic interfaces whose geometry and compliance may be
distant from the tactile sensor. Dynamics models, contact-mode controllers, and
contact fields address this setting~\cite{shirai2023tactiletool,vandermerwe2022compliant,oller2024extrinsic,higuera2023ncf}.
It tests whether local hand sensing can identify remote tool contact, but not
necessarily whether two hands coordinate.

\paragraph{Two hands with a tool and workpiece.}
Here one hand actuates a tool and the other stabilizes or repositions the
workpiece. Brushing a plate, wiping an object, scrubbing, cutting, peeling, or
inserting against a held fixture produces a multi-hop path such as
$H_L\!\rightarrow T\!\rightarrow W\!\rightarrow H_R$. BUDS establishes the
importance of stabilizer/actor role assignment~\cite{grannen2023stabilize}; TACO
documents broad bimanual tool--action--object semantics~\cite{liu2024taco}; and
TAMEn includes tactile-aware data collection in tasks such as dish
washing~\cite{wu2026tamen}. Tool use is therefore an established area rather than
an uncrowded task category. The narrower mapped cell---deployment-observable
cross-hand load propagation with a held workpiece and matched representation
alternatives---contains comparatively few direct studies.

\subsection{What a useful representation must retain}

Across topologies, five information requirements recur.

\begin{enumerate}
  \item \textbf{Interface identity and quality.} A representation must preserve
  which finger/link touches which entity, contact confidence, sensor quality,
  saturation, and missing-data masks. Treating a damaged taxel as zero pressure
  confounds absence of contact with absence of measurement.

  \item \textbf{Common spatial meaning.} Local tactile features require a link,
  hand, object, or task frame before contacts on separate bodies can be compared.
  Spatial anchoring~\cite{huang2025sata} and object-centric contact fields
  ~\cite{higuera2023ncf,ma2026semanticcontact} are different solutions to this
  frame-assignment problem.

  \item \textbf{Temporal event structure.} Contact onset, load acceptance,
  incipient slip, release, impact, and recontact occur on different time scales.
  A pooled window may recognize a state while obscuring its transition.

  \item \textbf{Cross-interface dependence.} The representation should allow a
  downstream model to distinguish coincident local events from a mechanically
  mediated relation. For handover, receiver loading should predict giver
  unloading and release readiness; for tool use, tool--workpiece changes should
  affect the stabilizing hand.

  \item \textbf{Uncertainty and observability.} Contact edges and remote loads are
  latent. A credible model should express uncertainty when sensors are occluded,
  saturated, delayed, or out of distribution instead of returning an
  overconfident topology.
\end{enumerate}

\begin{table}[t]
\centering
\caption{Topology-specific representation questions. ``Direct'' means that an
interface is measured locally; ``remote'' denotes an interface whose effects must
be inferred through another entity.}
\label{tab:topology_requirements}
\begin{tabularx}{\linewidth}{P{0.22\linewidth} P{0.20\linewidth} Y Y}
\toprule
\textbf{Topology} & \textbf{Critical transition} & \textbf{Observable evidence} & \textbf{Core representation test} \\
\midrule
Multi-finger in-hand & stick--slip / regrasp & several direct finger contacts & Can local evidence be aligned through kinematics and object motion? \\
Shared-object bimanual & unilateral--bilateral support & two hand streams, proprioception, optional vision & Is cross-hand dependence captured beyond concatenation? \\
Quasi-static handover & giver--dual--receiver & load redistribution and contact onsets & Can phase, load share, and safe release be predicted without privileged contacts? \\
Dynamic handover & release--flight--impact & pre-release contact and post-impact impulse & Does the representation anticipate impact and recovery under temporal discontinuity? \\
Single-hand tool use & free--working contact & grasp touch plus remote tool response & Can extrinsic tool contact be localized and its mode inferred? \\
Cooperative tool use & support/actor coupling changes & direct hand contacts; remote tool--workpiece effects & Does structure transfer across a multi-hop, role-asymmetric topology? \\
\bottomrule
\end{tabularx}
\end{table}

This topology view also clarifies what is \emph{not} required. A model need not
reconstruct every Newton of the complete scene wrench to support safe control;
relative load share, future interface wrench, slip risk, or a calibrated contact
phase may be sufficient. Conversely, high task success alone does not show that
the learned latent encodes any of these quantities. Representation claims need
diagnostic evidence.
